\section{Modelos implementados}

\subsection{MLP}
Se implementó el perceptrón multicapa de las diapositivas como una subclase de \texttt{nn.Sequential}. La entrada
de tamaño $(1, 28, 28)$ se aplana con \texttt{nn.Flatten} y atraviesa dos capas ocultas de 200 neuronas con activación
ReLU, seguidas de una capa de salida de 10 neuronas:
\[
784 \rightarrow 200 \rightarrow 200 \rightarrow 10
\]
La última capa devuelve directamente los \textit{logits}, sin aplicar \textit{softmax}, ya que la función de pérdida
de PyTorch (\texttt{cross\_entropy}) la incorpora internamente. El número de neuronas ocultas se lee del parámetro
\texttt{hidden} de la configuración, y con el valor por defecto el modelo tiene 199.210 parámetros.

\subsection{CNN}
Como punto de partida se tomó la red convolucional del artículo original de FedAvg~\cite{mcmahan}, formada por dos
convoluciones $5 \times 5$ de 32 y 64 canales, cada una seguida de ReLU y \textit{max pooling} $2 \times 2$, una
capa densa de 512 neuronas y la capa de salida. Esta red tiene 1.663.370 parámetros, de los que el 96.6\,\%
pertenecen a la capa densa, por lo que se eliminó dicha capa para cumplir con el límite de 100.000 parámetros,
manteniendo las convoluciones tal y como aparecen en el artículo. Además, se tomaron las siguientes decisiones de
diseño, orientadas al entrenamiento federado con datos no-IID:

\begin{itemize}
    \item \textbf{Sin BatchNorm}: esta capa normaliza con la media y la varianza de cada mini-lote, que con datos
    no-IID dependen de la distribución local del cliente y no se parecen a las del conjunto global, lo que degrada
    FedAvg~\cite{hsieh}.
    \item \textbf{Red ancha y poco profunda}: con datos no-IID, aumentar la profundidad empeora los resultados,
    mientras que aumentar la anchura los mejora~\cite{divergencia}. Por ello, los parámetros se emplean en canales
    y no en capas.
    \item \textbf{Aplanado y capa lineal}: en lugar de \textit{global average pooling}, que con datos no-IID deja
    una entrada demasiado pequeña al clasificador y empeora el resultado~\cite{divergencia}.
    \item \textbf{Sin \textit{dropout} ni \textit{weight decay}}, y con los mismos hiperparámetros de entrenamiento
    que el MLP, de modo que la comparación entre ambos modelos solo dependa de la arquitectura.
\end{itemize}

A partir de este diseño se consideraron tres variantes:

\begin{itemize}
    \item \textbf{CNN sin GroupNorm} (83.466 parámetros): la red descrita, que fue la variante inicial.
    \item \textbf{CNN $3 \times 3$ sin GroupNorm} (50.186 parámetros): la misma red con filtros $3 \times 3$, con
    menor campo receptivo y menor coste.
    \item \textbf{CNN con GroupNorm} (83.658 parámetros): la variante inicial con GroupNorm~\cite{groupnorm} tras
    cada convolución, con 8 grupos. A diferencia de BatchNorm, normaliza cada ejemplo por separado, sin emplear
    estadísticas del mini-lote, por lo que no presenta el problema anterior y recupera gran parte de la pérdida
    que sufren los modelos normalizados con datos no-IID~\cite{hsieh}.
\end{itemize}

Para la elección de la mejor variante de CNN se tomó como criterio la \textit{accuracy} final sobre el conjunto de
test, calculada como la media de las 5 últimas rondas de cada ejecución y promediada entre 3 semillas, con IID y
Dirichlet $\alpha=0.1$. Se mantendría la variante inicial salvo que otra la superase en $\alpha=0.1$ por más de una
desviación típica:

\begin{table}[H]
    \centering
    \caption{\textit{Accuracy} final de las variantes de CNN (media $\pm$ desviación típica entre 3 semillas)}
    \label{tab:cnn_variantes}
    \setlength{\tabcolsep}{5pt}
    \begin{tabular}{lccc}
        \hline
        Variante & Parámetros & IID & Dirichlet $\alpha=0.1$ \\
        \hline
        CNN sin GroupNorm & 83.466 & $0.883 \pm 0.002$ & $0.782 \pm 0.046$ \\
        CNN $3 \times 3$ sin GroupNorm & 50.186 & $0.880 \pm 0.002$ & $0.775 \pm 0.037$ \\
        CNN con GroupNorm & 83.658 & $0.903 \pm 0.003$ & $0.833 \pm 0.033$ \\
        \hline
    \end{tabular}
\end{table}

Emplear los datos de test para una decisión de diseño no resulta adecuado, ya que este conjunto debe ser
independiente de cualquier decisión tomada sobre el modelo. Sin embargo, la implementación de un conjunto de
validación para realizar la selección escapa al alcance de esta práctica, por lo que se decidió trabajar de esta
manera.
\\

La variante $3 \times 3$ queda dentro de la variación entre semillas, mientras que la variante con GroupNorm supera
a la inicial en 2.0 puntos con IID y en 5.2 con $\alpha=0.1$, por encima de la desviación típica de ambas. Por ello,
se eligió la CNN con GroupNorm, implementada en la clase \texttt{CNN}.
